% !TeX root = ../papers/nc-p3.tex
\section{滑らかな六次曲線と孤立残余の全排除}
\label{sec:smooth-sextic-exclusion}
\begin{theorem}\label{thm:smooth-sextic-exclusion}
標準有限Frobenius候補で四辺の左右の射影ファイバーが空か一点なら、
点グラフが滑らかな整六次曲線 $C$ とそれに交わらない零次元スキームの直和になることはない。
\end{theorem}
\begin{proof}[種数の限定]
各位置の空間次数は6で一定である。平面六次曲線なら種数は10であり、
非退化な空間六次曲線のCastelnuovo上界4により全位置で平面に入る。
両平面の積の中で全関係の基点は次元1以下なので、一般の三線型結合は三本の
双一次式による完全交叉曲線を定める。その両側次数は3であり、六次曲線を含めない。
従って全位置で非退化である。系\ref{cor:point-graph-excess}で孤立残余の長さは $2g-6$、
Castelnuovoにより $g\le4$ である。
$g=4$ なら $L_i$ の四つの独立切断とRiemann--Rochにより次数零の
$\omega_C/L_i$ が自明であり、全位置が同じ標準埋込みになる。
$h^0(\omega_C^2)=9$ なのでこれを含む二次曲面がある。
独立な二式は存在しない。完全交叉なら次数4に過ぎず、共通平面因子なら
整な非退化曲線が平面または直線に入ってしまう。
従って曲線上の双一次関係は全てこの二次曲面の対角上でも零であり、点グラフの次元に反する。
残る $g=3$ の場合は孤立部の長さが零で、非超楕円の場合は前節で排除した。
\end{proof}

\begin{proof}[超楕円な場合の分裂型]
二重被覆を $f:C\to\mathbb P^1$、$J=f^*\mathcal O(1)$、
$W=H^0(\mathbb P^1,\mathcal O(1))$ とする。$L_i$ は次数6の完全な四切断の偏極で、
$\omega_C=J^2$、$L_i/J^2$ は次数2で切断を持たない（持てば二点を分離できない）。
$f_*L_i$ は階数2、次数2であり、$f_*L_i(-2)$（直線上でのひねり）は切断を持たない。
従って分裂型は $\mathcal O(1)^2$ である。
$P_i=L_i/J$ は基点自由な次数4の鉛筆で、$U_i=H^0(P_i)$ とすると
$f_*P_i=U_i\otimes\mathcal O$、$V_i=W\otimes U_i$ となる。
各埋込みは二鉛筆の積で滑らかなSegre二次曲面に入る。
もし $P_i=P_{i+1}$ なら二次曲面の対角も点グラフに入り矛盾する。
ここでも六次の非退化整曲線を含む二次曲面は一つだけである。

各二文字で積の束写像の核を $K_i$ と書く。
\[
 0\longrightarrow K_i\longrightarrow
 U_i\otimes U_{i+1}\otimes\mathcal O\longrightarrow f_*(P_i P_{i+1})\longrightarrow0 .
\]
両束 $P_i$ は有限平坦二重被覆の各ファイバー上でその全切断空間を $U_i$ から受け取るため、
局所Artin環上の可逆加群を自明化した通常の乗法により最後の射は全射である。
従って $K_i$ は階数2、次数 $-4$ であり、$H^0(K_i)=0$ である。
最後の消滅は $P_i$ の鉛筆列を $P_{i+1}$ でひねたときの次数零の核束
$P_{i+1}/P_i$ が非自明なためである。
直和分解の可能性は $(-1,-3),(-2,-2)$ だけであり、いずれも $H^1(K_i(2))=0$。
よって $V_i\otimes V_{i+1}$ から曲線の全二次切断への積は全射、六関係は全ての核である。
このうち四本は $\bigwedge^2 W\otimes U_i\otimes U_{i+1}$、残り二本は
$H^0(K_i(2))$ で与えられる。

分裂型 $(-1,-3)$ なら最後の二式は各基点 $z\in\mathbb P^1$ で一形式の倍にしかならない。
両側の第一鉛筆を $z$ に固定した $\mathbb P(U_i^*)\times\mathbb P(U_{i+1}^*)$ では
四つの交代式は全て零で、残りが正次元の零点を持つ。
実際の点グラフはこの基点上で有限なので矛盾する。従って全て $K_i=\mathcal O(-2)^2$ である。

三文字の全積核の束を調べる。二つの部分束 $K_i\otimes U_{i+2}$ と $U_i\otimes K_{i+1}$ を
階数8の自明束に入れ、その交叉を $B_i$ とする。和の余核は $f_*(P_iP_{i+1}P_{i+2})$ である。
実際、長さ2の各ファイバーの二次元環 $D$ で、三つの可逆加群を自明化すると
$D^{\otimes3}$ を隣接二箇所の乗法核で割ったものは $D$ である。
単位を挿入することでどの三テンソルも全積と二つの単位に等しくなるためである。
従って和は定数階数、$B_i$ は階数2の部分束で、
$\deg B_i=-8-8+(12-4)=-8$ となる。
$B_i(3)$ の次数は $-2$ で、$K_i(3)\otimes U_{i+2}=\mathcal O(1)^4$ に入るため、
二つの直和因子の次数は高々1である。よって $h^0(B_i(3))\le2$。

四つの交代式を各辺で割った三文字空間は
$\operatorname{Sym}^3 W\otimes U_i\otimes U_{i+1}\otimes U_{i+2}$ で次元32である。
$W$ の隣り合う二文字を自由に交換し、三つの $U$ の係数は独立に保持するためである。
残る左と右の二次関係の挿入はそれぞれ $H^0(K_i(3))\otimes U_{i+2}$,
$U_i\otimes H^0(K_{i+1}(3))$ の八次元空間を与える。
$K_i(2)$ は自明束なのでここで $W$ との積は全切断を張る。
両者の交叉は $H^0(B_i(3))$ である。従って三文字商の次元は高々
$32-8-8+2=18$、有限Frobenius候補に必要な20に反する。
\end{proof}
